% !TEX program = lualatex
% Transparencias de Fundamentos de las Matemáticas I
% Clase 14: Anillos y cuerpos

\documentclass[aspectratio=169,11pt]{beamer}

\usepackage{fontspec}
\usepackage[spanish,es-nodecimaldot]{babel}
\usepackage{amsmath,amssymb,mathtools}
\usepackage{booktabs}
\usepackage{csquotes}
\usepackage{xcolor}

\usetheme{Madrid}
\usecolortheme{default}
\setbeamertemplate{navigation symbols}{}
\setbeamertemplate{footline}[frame number]

\definecolor{FdMBlue}{RGB}{20,55,100}
\definecolor{FdMRed}{RGB}{130,30,30}
\definecolor{FdMGreen}{RGB}{25,100,60}

\setbeamercolor{structure}{fg=FdMBlue}
\setbeamercolor{title}{fg=white,bg=FdMBlue}
\setbeamercolor{frametitle}{fg=white,bg=FdMBlue}
\setbeamercolor{block title}{fg=white,bg=FdMBlue}
\setbeamercolor{block body}{bg=blue!3}
\setbeamercolor{alerted text}{fg=FdMRed}

\setmainfont{Latin Modern Roman}
\setsansfont{Latin Modern Sans}
\setmonofont{Latin Modern Mono}

\newcommand{\N}{\mathbb{N}}
\newcommand{\Z}{\mathbb{Z}}
\newcommand{\Q}{\mathbb{Q}}
\newcommand{\R}{\mathbb{R}}
\newcommand{\C}{\mathbb{C}}
\newcommand{\Pow}{\mathcal{P}}
\newcommand{\id}{\operatorname{id}}
\newcommand{\mcd}{\operatorname{mcd}}
\newcommand{\mcm}{\operatorname{mcm}}

\title[FdM I -- Clase 14]{Anillos y cuerpos}
\subtitle{Estructuras con dos operaciones}
\author{Antonio Falcó Montesinos}
\institute{Universidad CEU Cardenal Herrera}
\date{Fundamentos de las Matemáticas I}

\begin{document}

\begin{frame}
  \titlepage
\end{frame}

\begin{frame}{Objetivos de la clase}
\begin{itemize}
  \item Definir anillos y anillos con unidad.
  \item Definir cuerpos.
  \item Relacionar estas estructuras con sistemas numéricos conocidos.
\end{itemize}
\end{frame}

\begin{frame}{Mapa de la clase}
\tableofcontents
\end{frame}

\section{Anillos}

\begin{frame}{Anillos}
\begin{block}{Idea}
\begin{itemize}
  \item Un anillo tiene dos operaciones: suma y producto.
  \item La suma forma un grupo abeliano.
  \item El producto es asociativo y distributivo respecto de la suma.
\end{itemize}
\end{block}
\end{frame}

\begin{frame}{Anillos}
\begin{block}{Anillo con unidad}
\begin{itemize}
  \item Un anillo con unidad posee un neutro multiplicativo \(1\).
  \item En estas notas se dirá explícitamente cuando se exija unidad.
  \item \(\mathbb{Z}\) con suma y producto usuales es un anillo con unidad.
\end{itemize}
\end{block}
\end{frame}

\begin{frame}{Anillos}
\begin{block}{Distributividad}
\begin{itemize}
  \item La propiedad distributiva conecta las dos operaciones.
  \item \(a(b+c)=ab+ac\) y \((a+b)c=ac+bc\).
  \item Sin distributividad, no hay estructura de anillo.
\end{itemize}
\end{block}
\end{frame}

\section{Cuerpos}

\begin{frame}{Cuerpos}
\begin{block}{Definición}
\begin{itemize}
  \item Un cuerpo permite sumar, restar, multiplicar y dividir por elementos no nulos.
  \item \((K\setminus\{0\},\cdot,1)\) debe ser grupo abeliano.
  \item Además, el producto distribuye respecto de la suma.
\end{itemize}
\end{block}
\end{frame}

\begin{frame}{Cuerpos}
\begin{block}{Ejemplos}
\begin{itemize}
  \item \(\mathbb{Q}\), \(\mathbb{R}\) y \(\mathbb{C}\) son cuerpos.
  \item \(\mathbb{Z}\) no es cuerpo: \(2\) no tiene inverso multiplicativo entero.
  \item Las matrices cuadradas forman un anillo con unidad, pero no un cuerpo si el tamaño es al menos dos.
\end{itemize}
\end{block}
\end{frame}

\begin{frame}{Cuerpos}
\begin{block}{Mapa conceptual}
\begin{itemize}
  \item Cuerpo \(\Rightarrow\) anillo con unidad.
  \item Anillo con unidad no implica cuerpo.
  \item La diferencia clave está en los inversos multiplicativos de los no nulos.
\end{itemize}
\end{block}
\end{frame}

\section{Profundización}

\begin{frame}{Profundización}
\begin{block}{Dos operaciones, dos papeles}
\begin{itemize}
  \item La suma organiza la estructura aditiva.
  \item El producto introduce multiplicación y distributividad.
  \item La interacción entre ambas operaciones es lo que caracteriza a un anillo.
\end{itemize}
\end{block}
\end{frame}

\begin{frame}{Profundización}
\begin{block}{Por qué \(\mathbb{Z}\) no es cuerpo}
\begin{itemize}
  \item En un cuerpo, todo elemento no nulo debe tener inverso multiplicativo.
  \item En \(\mathbb{Z}\), no existe \(b\in\mathbb{Z}\) tal que \(2b=1\).
  \item Por tanto, \(\mathbb{Z}\) es anillo con unidad, pero no cuerpo.
\end{itemize}
\end{block}
\end{frame}

\begin{frame}{Profundización}
\begin{block}{Lectura hacia adelante}
\begin{itemize}
  \item Los racionales se construirán para permitir división exacta entre enteros.
  \item Los reales completarán a los racionales.
  \item Los complejos permitirán resolver ecuaciones sin solución real.
\end{itemize}
\end{block}
\end{frame}

\begin{frame}{Cierre}
\begin{itemize}
  \item Los sistemas numéricos se entienden mejor como estructuras.
  \item La suma y el producto cumplen papeles distintos.
  \item La próxima clase inicia números naturales e inducción.
\end{itemize}
\end{frame}

\begin{frame}{Trabajo recomendado}
\begin{itemize}
  \item Releer las secciones correspondientes del manual.
  \item Identificar definiciones, símbolos nuevos y enunciados principales.
  \item Preparar dudas y ejemplos para el seminario asociado.
\end{itemize}
\end{frame}

\end{document}
